\chapter[面向对象编程]{面向对象编程（OOP）}
\label{ch:05-outline}

设计类时，先确定对象需要保存的状态和提供的操作，再考虑继承关系。实例属性、类属性与协作对象各有职责，不能只凭语法相似就混用。


\begin{figure}[htbp]
  \centering
  \begin{tikzpicture}[
    book mindmap,
    level 1 concept/.append style={minimum size=22mm,text width=21mm}
  ]
    \node[text=white,minimum size=25mm,text width=19mm] {第 5 章}
      child[grow=163,level distance=38mm,concept color=DeepBlue] {
        node[text=white] {类与实例}
      }
      child[grow=17,level distance=38mm,concept color=PythonYellow] {
        node {属性与封装}
      }
      child[grow=-90,level distance=28mm,concept color=PythonBlue!22!white] {
        node {继承与共享状态}
      };
  \end{tikzpicture}
  \caption*{实例状态、属性校验与继承协作。}
\end{figure}


\section{类与对象}

类像模具，实例是做出的某一块饼干。在 \code{__init__} 里设置实例属性。
\code{__str__} 用于日常显示，\code{__repr__} 为调试提供对象的可辨识表示。

方法里的 \code{self} 就是这一次调用对应的实例。它让同一段方法代码操作各自的状态。
下面用一条联系人记录对照两种显示，随后用同步任务检查实例属性。


\begin{lstlisting}[title={\code{contact_repr.py}}]
class Contact:
    def __init__(self, name):
        self.name = name

    def __str__(self):
        # 日常显示用简洁名字。
        return self.name

    def __repr__(self):
        # 调试时保留类型名和字符串边界。
        return f"Contact(name={self.name!r})"


person = Contact("小雪")
print(person)
print(repr(person))
\end{lstlisting}


\textbf{运行核对}\enspace 预期输出：
\begin{lstlisting}[language={},numbers=none]
小雪
Contact(name='小雪')
\end{lstlisting}

\par


\section{封装与 @property}

将 property 当成属性前面的一道门。调用处照常写 \code{job.name = value}，
门内统一处理去空格和校验。访问属性时可以通过 getter 返回整理后的值，
赋值时则经过 setter。属性内部用另一个名字存储，避免再次调用同一个 setter。

单下划线是内部使用的约定。它方便读代码的人识别边界，具体校验仍由代码执行。
在 PyCharm 中先找属性的使用处，再检查改名或添加校验会影响哪些调用。

\subsection*{动手｜让两个实例各管各的状态}

用两个同步任务检查属性放在哪里。任务名经过同一个属性入口清理，日志列表在初始化时创建。这样只给第一个任务记一条消息，第二个任务的列表就会保持为空。

\begin{lstlisting}[title={\code{sync_jobs.py}}]
class SyncJob:
    def __init__(self, name):
        self.name = name
        # 每个实例在初始化时建自己的日志。
        self.log = []

    @property
    def name(self):
        return self._name

    @name.setter
    def name(self, value):
        cleaned = value.strip()
        if not cleaned:
            raise ValueError("任务名不能为空")
        self._name = cleaned


# 给一个实例写日志，再检查另一个实例。
first = SyncJob(" notes.md ")
second = SyncJob("todo.md")
first.log.append("queued")
print(first.name, first.log)
print(second.name, second.log)
\end{lstlisting}

\textbf{运行核对}\enspace 正常输出两行，notes.md 后面有一条日志，todo.md 后面是空列表。把任务名换成只含空格的字符串，可以复现校验失败。

\section{继承与常见设计误区}

类属性像挂在模具旁边的一张共享纸条。实例属性则写在每块饼干自己的标签上。
若把日志列表写在类体里，所有实例都可能看到同一份列表。
上节把列表放进 \code{__init__}，就让每次创建对象时各建一份。

继承时，先找真正共用的行为。\code{super()} 按方法解析顺序把调用交给后面的类。
多继承的协作要求各类保持能衔接的参数和调用约定，先用单继承把路径看清。


\begin{lstlisting}[title={\code{inherit_name.py}}]
class Job:
    def label(self):
        return "sync"


class NamedJob(Job):
    def label(self):
        # super 取父层行为，再补当前层的信息。
        return super().label() + ":notes.md"


print(NamedJob().label())
print([cls.__name__ for cls in NamedJob.__mro__])
\end{lstlisting}


\textbf{运行核对}\enspace 预期输出：
\begin{lstlisting}[language={},numbers=none]
sync:notes.md
['NamedJob', 'Job', 'object']
\end{lstlisting}

\par

\begin{myunderstanding}
实例属性给每个对象自己的状态，类属性适合真正共享的配置。继承共用行为时，要让每层的职责和参数能接起来。
\end{myunderstanding}

\begin{pycharmpractice}
在属性 setter 设断点，给实例赋值后观察输入和内部存储字段。查找用法可以帮助核对属性更名的影响。
\end{pycharmpractice}

\begin{pitfall}
在 setter 里再次赋值给同一个属性，会递归调用自己。用另一个内部字段保存校验后的值。
\end{pitfall}

\subsection*{进阶边界演练}
给类增加字段前，先问这个字段属于每个实例，还是整类共享。需要组合另一种能力时，也可以让实例持有一个协作对象。遇到多继承，再查看 MRO，把每一次 super 的去向画出来。
